\chapter{Karyotype (Chromosome Analysis)}

\section*{What Is This Test?}
A karyotype is the complete set of chromosomes in an organism, analyzed by photographing metaphase chromosomes and arranging them in a standardized format (the karyogram) according to chromosome size, centromere position, and banding pattern. The test involves the culture of living cells (peripheral blood lymphocytes, bone marrow, amniotic fluid cells, chorionic villus tissue, or fibroblasts) to stimulate cell division, arresting cells in metaphase with colchicine (a microtubule inhibitor), releasing the chromosomes by hypotonic lysis, fixing them on glass slides, and staining with Giemsa to produce a G-banded karyotype with approximately 400--550 bands per haploid set (high-resolution banding $>$800 bands). Each chromosome is identified by its number (1--22 autosomes plus X and Y sex chromosomes), and structural abnormalities are described using the International System for Human Cytogenomic Nomenclature (ISCN) \cite{mcgowanjordan2020iscn}. Karyotype analysis can detect aneuploidy (loss or gain of whole chromosomes, e.g., trisomy 21 in Down syndrome, monosomy X in Turner syndrome), large structural rearrangements (translocations, deletions, duplications, inversions, ring chromosomes, isochromosomes), and marker chromosomes of unknown origin. Karyotyping is essential for constitutional genetic disorders (prenatal diagnosis, pediatric developmental delay, recurrent pregnancy loss) and hematologic malignancies (leukemia, myelodysplastic syndromes), where clonal chromosomal abnormalities provide diagnostic, prognostic, and therapeutic guidance.

\section*{Alternative Names}
Chromosome Analysis; Cytogenetic Analysis; Karyogram; G-Banding; Metaphase Chromosome Analysis; Giemsa Banding; Constitutional Karyotype; Leukemia Karyotype

\section*{Principle}
Living cells are cultured in vitro with a mitogen: phytohemagglutinin (PHA) for T lymphocytes (constitutional karyotype from peripheral blood), pokeweed mitogen or IL-2 lipopolysaccharide for B lymphocytes, or spontaneous proliferation for bone marrow aspirate (leukemia karyotype). Cells are incubated for 48--72 hours (peripheral blood) or 24--48 hours (bone marrow) at 37 degrees C in a 5\% CO2 atmosphere. Colchicine (colcemid) is added to arrest cells in metaphase by disrupting microtubule assembly of the mitotic spindle. Cells are swollen in hypotonic potassium chloride solution, which separates chromosome spreads but maintains nuclear integrity. The cells are fixed in Carnoy fixative (3:1 methanol:glacial acetic acid), dropped onto glass slides, and air-dried. Slides are treated with trypsin to partially digest chromosomal proteins and then stained with Giemsa stain, which binds differentially to AT-rich (dark G-bands) and GC-rich (light G-bands) DNA regions, producing a characteristic banding pattern. At least 20 metaphase cells are analyzed, and a minimum of 5 cells are fully karyotyped. The cytogeneticist identifies numerical and structural chromosome abnormalities. The resolution is limited to approximately 5--10 megabases at the standard 400--550 band level; deletions or duplications smaller than 5--10 Mb are not detectable by karyotype.

\section*{History}
The correct human chromosome number (46) was established in 1956 by Joe Hin Tjio and Albert Levan \cite{tjio1956chromosome}. The first chromosomal abnormality associated with a human disease was the discovery of trisomy 21 in Down syndrome by Jerome Lejeune in 1959 \cite{lejeune1959etude}. The Philadelphia chromosome was discovered in 1960 by Nowell and Hungerford \cite{nowell1960minute}. Chromosome banding techniques (Q-banding by Caspersson, G-banding by Seabright) were developed in the early 1970s, enabling identification of individual chromosomes and precise breakpoint mapping \cite{caspersson1970identification, seabright1971rapid}. Fluorescence in situ hybridization (FISH) was developed in the 1980s as a complementary technique \cite{pinkel1986cytogenetic}. The International System for Human Cytogenomic Nomenclature (ISCN) was first published in 1960 and is updated every 2--4 years \cite{mcgowanjordan2020iscn}. Chromosomal microarray analysis (CMA, array CGH) and next-generation sequencing are gradually complementing and, in many settings (particularly pediatric developmental delay and autism spectrum disorder), supplanting karyotype due to their much higher resolution \cite{miller2010consensus, south2013acmg}.

\section*{Why This Test Is Ordered}
Constitutional karyotype (peripheral blood): Prenatal diagnosis (amniocentesis, chorionic villus sampling) for advanced maternal age, abnormal prenatal screening, or ultrasound anomalies. Evaluation of unexplained developmental delay, intellectual disability, autism spectrum disorder, or multiple congenital anomalies in a child. Evaluation of recurrent pregnancy loss ($\geq$2 consecutive miscarriages). Evaluation of primary amenorrhea, primary ovarian insufficiency, or suspected Turner syndrome. Evaluation of infertility in men with azoospermia or severe oligospermia (Klinefelter syndrome 47,XXY). Confirmation of suspected trisomy 21, 18, or 13 in a newborn with dysmorphic features. Hematologic malignancy karyotype (bone marrow): diagnosis, risk stratification, and monitoring of acute leukemias (AML, ALL), myelodysplastic syndromes (MDS), myeloproliferative neoplasms (MPN), chronic myeloid leukemia (CML), and lymphoma.

\section*{Is This a Primary or Secondary Test}
For constitutional disorders, karyotype is a primary diagnostic test but is increasingly being replaced by chromosomal microarray (CMA) for pediatric developmental delay and autism \cite{miller2010consensus}. For hematologic malignancies, karyotype at diagnosis is an essential primary test required for risk stratification by WHO, NCCN, and ELN criteria \cite{dohner2017diagnosis, swerdlow2017who}.

\section*{How This Test Is Performed}
Peripheral blood: a blood sample is collected in a sodium heparin (green-top) tube. Bone marrow: 1--2 mL of bone marrow aspirate collected in a sodium heparin tube (the first aspirate is preferred for cytogenetics). Amniotic fluid: 20--30 mL collected by amniocentesis. Chorionic villus tissue: obtained by CVS. Products of conception: placental and fetal tissue from miscarriage. The specimen must be received by the cytogenetics laboratory within 24 hours (peripheral blood) or transported immediately (bone marrow), and it must be sterilely collected to prevent microbial contamination of cultures.

\section*{What Preparation Is Needed}
No special preparation for peripheral blood collection. For bone marrow aspirate for hematologic malignancy karyotype, the bone marrow aspirate should ideally be collected before transfusion of blood products if possible, as transfused leukocytes may overgrow the culture and obscure the malignant karyotype.

\section*{Contraindications and Precautions}
The major limitation is resolution: deletions or duplications <5--10 Mb are not detectable by karyotype. Balanced rearrangements (translocations, inversions) may appear ``normal'' if the rearrangement is cryptic (too small to visualize) or if the breakpoints are at the cytogenetic band level. Maternal cell contamination (MCC) can occur in prenatal specimens (particularly CVS), resulting in a 46,XX maternal karyotype rather than the fetal karyotype. A failure to culture (no growth or poor growth of the specimen) occurs in approximately 0.5--2\% of peripheral blood karyotypes, 5--10\% of bone marrow karyotypes (due to low cellularity or previous chemotherapy), and <1\% of prenatal karyotypes. Bone marrow karyotypes in CLL require B-cell mitogens (CpG-ODN and IL-2) because CLL cells do not divide spontaneously. Heparin is the only acceptable anticoagulant; EDTA and citrate inhibit cell division and will result in culture failure. Specimen should be transported at room temperature (NOT refrigerated and NOT frozen).

\section*{Risks}
Minimal risk from venipuncture. For prenatal diagnosis (amniocentesis, CVS): procedure-related pregnancy loss risk approximately 0.1--0.3\% (1:300 to 1:1,000). For bone marrow aspiration: pain at the aspiration site, bleeding, infection (rare <1\%).

\section*{What Happens After the Test}
Results are available in: constitutional karyotype (PHA-stimulated peripheral blood) 7--10 days; prenatal karyotype (amniotic fluid, CVS) 7--14 days; bone marrow karyotype for leukemia 5--14 days (depending on whether the cells are dividing spontaneously or stimulated). The report includes the total number of cells analyzed, the karyotype written in ISCN nomenclature, and an interpretation by a certified cytogeneticist. Abnormal results should be discussed with a genetic counselor (constitutional) or hematologist-oncologist (malignancy).

\section*{Understanding Results}

\subsection*{Normal Karyotype}
Constitutional: 46,XX (normal female) or 46,XY (normal male). Bone marrow (hematologic malignancy): 46,XX or 46,XY (normal diploid karyotype). A normal karyotype in a leukemic patient (cytogenetically normal AML, CN-AML) is found in approximately 40--50\% of AML and 30--40\% of ALL and carries an intermediate-risk prognosis (with additional molecular markers such as NPM1, FLT3-ITD, and CEBPA further refining risk) \cite{dohner2017diagnosis}.

\subsection*{Abnormal Karyotype}
Constitutional aneuploidy: Trisomy 21 (47,XX,+21 or 47,XY,+21): Down syndrome. Trisomy 18 (47,XX,+18 or 47,XY,+18): Edwards syndrome. Trisomy 13 (47,XX,+13): Patau syndrome. Monosomy X (45,X): Turner syndrome. 47,XXY: Klinefelter syndrome. 47,XYY: Jacob syndrome. Constitutional structural rearrangements: Balanced reciprocal translocations (e.g., 46,XX,t(11;22)(q23;q11)): carriers are usually phenotypically normal but have an increased risk of infertility, recurrent pregnancy loss, and offspring with unbalanced chromosomal abnormalities. Robertsonian translocations (e.g., 45,XY,der(13;14)(q10;q10)): the most common chromosomal rearrangement in humans. Hematologic malignancy: t(9;22)(q34;q11) = Philadelphia chromosome (CML, Ph+ ALL). t(15;17)(q24;q21) = PML-RARA (acute promyelocytic leukemia, APL). t(8;21)(q22;q22) = RUNX1-RUNX1T1 (AML, favorable risk). inv(16)(p13.1q22) = CBFB-MYH11 (AML, favorable risk). 11q23/MLL (KMT2A) rearrangements (AML and ALL, intermediate to high risk). Complex karyotype ($\geq$3 chromosomal abnormalities): high risk in AML, MDS. Monosomal karyotype: very high risk in AML. Deletion 5q, deletion 7q/monosomy 7: high risk in MDS and AML. Deletion 17p: high risk (TP53 loss) in multiple myeloma, CLL.

\section*{Normal Results}
Constitutional karyotype: 46,XX (female) or 46,XY (male) in $\geq$20 metaphase cells with no clonal abnormalities. Bone marrow karyotype: normal diploid karyotype 46,XX or 46,XY in $\geq$20 metaphases. A ``normal'' result does not exclude submicroscopic deletions, duplications, or point mutations --- these require CMA or molecular genetic testing. Results are reported as: 46,XX (normal female karyotype); 46,XY (normal male karyotype).

\section*{Follow-up Tests}
Chromosomal microarray analysis (CMA/array CGH): provides genome-wide copy number analysis at a resolution of 20--100 kb (100x higher than karyotype) and can detect submicroscopic deletions and duplications \cite{miller2010consensus, south2013acmg}. FISH (fluorescence in situ hybridization) using locus-specific probes to confirm a suspected microdeletion syndrome (22q11.2, 15q11.2, 7q11.23). MLPA (multiplex ligation-dependent probe amplification) for subtelomeric rearrangements and targeted gene dosage analysis. NGS-based targeted gene panels or whole-exome sequencing for monogenic disorders that are not detectable by karyotype. Molecular genetic testing for NPM1, FLT3-ITD/TKD, CEBPA, IDH1/IDH2, RUNX1, TP53 and other mutations in cytogenetically normal AML for risk stratification. Prenatal karyotype with concurrent CMA for fetuses with structural anomalies on ultrasound.

\section*{References}
\bibliographystyle{unsrtnat}
\bibliography{references}
\clearpage
